\documentclass{article}

% Anonymous submission. For the camera-ready version, replace the line below
% with \usepackage[final]{neurips_2026} and uncomment the author block.
\usepackage{neurips_2026}

\usepackage[utf8]{inputenc}
\usepackage[T1]{fontenc}
\usepackage{hyperref}
\usepackage{url}
\usepackage{booktabs}
\usepackage{amsfonts}
\usepackage{amsmath}
\usepackage{nicefrac}

\usepackage{graphicx}
\usepackage{xcolor}

\title{An Absolute Floor for Forecasting, and What It Reveals About Temporal Foundation Models}

\author{%
  Anonymous Author(s)\\
  Affiliation\\
  Address\\
  \texttt{email}
}

% Camera-ready author block.
% \author{%
%   Sarim Chaudhry\\
%   Purdue University\\
%   \texttt{chaud158@purdue.edu}
% }

\begin{document}

\maketitle

\begin{abstract}
Every metric in common use for forecasting is relative. Scaled errors and probabilistic scores report that one model beat another, and none reports how far either sits from what is achievable, because on collected data nobody knows what is achievable. We remove that limitation on a family of systems where the optimum is computable. A chaotic system produces information at a rate equal to the sum of its positive Lyapunov exponents, and that rate is the minimum a tracker must be supplied to stay locked onto a trajectory. We build an evaluation around it. A predictor runs on its own output and receives a correction only when it drifts past a chosen precision, and the correction stream in bits per unit time is the score. Sweeping the precision separates two causes of drift that error metrics conflate, since the cost of following a deterministic source does not change with precision while the cost of following a model's own mistakes does. On five systems spanning entropy rates from zero to two nats per unit time, four pretrained time series foundation models across a twentyfold range in parameters all spend more information than repeating the last observed value, all show the precision response of a predictor with no internal dynamics, and improve with scale too slowly to reach the floor below $10^{16.8}$ parameters. A delay embedding with a few hundred effective parameters reaches the floor exactly where the floor is zero, so the gap belongs to the models rather than to the instrument.
\end{abstract}

\section{Introduction}

Suppose a forecaster reaches a held out error of $0.003$ on some series. Is that good? Nobody can say. The number means something only beside another number from another model on the same data, which gives an ordering rather than a distance, and cannot say whether the remaining error is the irreducible unpredictability of the process or a deficiency a better model would remove. Forecasting evaluation is relative in this way throughout \citep{makridakis2020m4, makridakis2022m5, godahewa2021monash, aksu2024gifteval}.

This matters more than it used to. Time series foundation models are trained once on large heterogeneous corpora and then applied to series they have never seen \citep{ansari2024chronos, woo2024moirai, das2024timesfm, rasul2023lagllama, goswami2024moment, shi2025timemoe}, and the claim attached to them is not that they fit one dataset well but that they have learned something general about how quantities evolve in time. A relative leaderboard cannot test that claim, since it compares models on series whose difficulty is unknown and which may overlap with pretraining data in ways that are hard to rule out \citep{sainz2023contamination, magar2022contamination}.

There is a family of systems where the difficulty is known. A deterministic chaotic system generates information at a definite rate. Nearby trajectories separate at a speed set by the positive Lyapunov exponents, and by Pesin's identity their sum equals the Kolmogorov and Sinai entropy \citep{pesin1977, eckmann1985}. That quantity has a direct operational reading. To keep a model locked onto the true trajectory at any fixed precision you must supply new information at exactly that rate, and the requirement does not soften as the precision tightens \citep{shaw1981, ruelle1989}. It is a floor, it follows from the equations, and no predictor can go below it.

We build an evaluation around this. The instrument is a predictive coder. After a short warm up the predictor never sees the true signal again, running forward on its own reconstructions, and whenever the reconstruction drifts more than a chosen precision from reality the coder emits a correction that snaps it back. We measure the correction stream in bits per unit time and subtract the floor, and what remains is excess information, an absolute score that needs no reference model to interpret. Sweeping the precision turns this into something sharper. Following a deterministic source costs the same bits per unit time however tightly you track it, because a finer precision means a smaller starting error that takes proportionally longer to grow by the same factor. A model's own mistakes do not behave that way, since their size is fixed by the model rather than by the threshold. So a predictor whose state genuinely tracks the system gives a flat response to the sweep and one matching surface patterns gives a rising one, and the slope separates the two cases with a single number.

Applying this to four pretrained models across five systems, we find that all of them spend more information than a trivial baseline, that their precision response is indistinguishable from having no internal model at all, and that scale reduces the excess too slowly to reach the floor below $10^{16.8}$ parameters.

\section{Related work}

Scaled and probabilistic scores, the evaluations built on them, and the suites comparing time series foundation models with the architectures preceding them are all relative, as is the work on how those models scale \citep{hyndman2006, gneiting2007, makridakis2020m4, godahewa2021monash, aksu2024gifteval, ansari2024chronos, woo2024moirai, das2024timesfm, kaplan2020, edwards2024scaling}.

Chaotic systems have served as forecasting benchmarks before \citep{gilpin2021chaos, gilpin2023scale}, and a striking observation from that work is that the largest Lyapunov exponent does not predict how well learned methods forecast different systems. Our framework suggests why, since the quantity governing a tracker's information cost is the full entropy rather than the leading exponent. Classical nonlinear analysis supplies our baselines \citep{takens1981, farmer1987, kantz2004} and the coding machinery is standard \citep{shannon1948, jayant1984, krichevsky1981}. Reading entropy as the rate at which a chaotic system must be re-informed goes back to early work on information flow \citep{shaw1981, grassberger1986}, compression as a measure of modelling quality has returned in the language setting \citep{deletang2024compression}, and related questions about minimal predictive memory come from computational mechanics \citep{crutchfield1989, shalizi2001}. Computing only when a prediction fails is familiar from event based sensing \citep{lichtsteiner2008, gallego2022}, and the self feeding regime we measure is where world models operate \citep{ha2018world, hafner2023dreamerv3, lecun2022path}.

\section{Method}

\paragraph{The coder.}
Let $x_t$ be the trajectory of a system observed at interval $\Delta t$ and standardised per channel. A predictor is any function that extends a context to a block of future values. The coder proceeds one step at a time, taking the predictor's value $\hat{x}_t$ and quantising its residual against the truth to a chosen precision $\delta$,
\begin{equation}
n_t = \operatorname{round}\!\left(\frac{x_t - \hat{x}_t}{\delta}\right),
\qquad
r_t = \hat{x}_t + n_t \delta ,
\end{equation}
where $n_t$ is the correction and $r_t$ is the reconstruction, which satisfies $|x_t - r_t| \le \delta/2$ by construction. The predictor is fed $r_t$ and never $x_t$, so the corrections form a genuine transmitted message rather than a diagnostic readout. This is the classical arrangement of differential quantisation \citep{jayant1984} with the linear predictor replaced by the model under test.

\paragraph{The score.}
The correction stream is sparse and correlated in time, so a marginal histogram overstates its cost and we estimate the rate with an adaptive context model using standard smoothing \citep{krichevsky1981}. Writing $B$ for the resulting bits per step, the score is the excess over the floor,
\begin{equation}
G(\delta) \;=\; \frac{B(\delta) \log 2}{\Delta t} \;-\; h_{\mathrm{KS}},
\qquad
h_{\mathrm{KS}} \;=\; \sum_{\lambda_i > 0} \lambda_i .
\end{equation}
Exponents come from the analytic Jacobian of each system \citep{benettin1980, wolf1985}, so the floor is exact to integrator precision rather than estimated from data. We repeat the measurement across a range of $\delta$ and fit $G$ against $\log_2(1/\delta)$, and the resulting slope is the diagnostic.

\paragraph{Event driven scheduling.}
A predictor doing its job should not need consulting at every step. The coder requests a block of future values, consumes it while no correction is needed, and asks for a new one only when a correction fires or the block runs out. Forward passes therefore scale with how often the model is wrong rather than with the length of the signal.

\paragraph{The benchmark.}
We use five continuous time systems spanning a wide range of entropy rates, namely van der Pol, Rossler, Halvorsen, Lorenz and Chen, with floors of $0.000$, $0.067$, $0.651$, $0.905$ and $2.061$ nats per unit time. Van der Pol is a limit cycle and its floor is exactly zero, which makes it the calibration case. Each is sampled at a fixed number of samples per Lyapunov time and standardised per channel, so the precision is in standard deviations and comparable across systems. Because the data is generated rather than collected, no pretrained model can have seen these trajectories in training. That is a proof rather than an audit \citep{sainz2023contamination, carlini2023memorization}.

\section{Analysis}
\label{sec:analysis}

Consider tracking a chaotic trajectory to within a tolerance $\delta$. Any error grows at a rate set by the positive exponents, so starting from the quantisation floor of order $\delta$ it reaches the tolerance again after a time of order $1/\lambda$, and the correction required then carries the information the system produced in that interval. The rate of injection is therefore $h_{\mathrm{KS}}$ and it does not depend on $\delta$, because a finer tolerance both shrinks the starting error and lengthens the time needed to grow by the same factor \citep{shaw1981, ruelle1989}.

That cancellation is the diagnostic. For a predictor whose state evolves with the system the residual after each correction is dominated by the quantisation floor, so the cancellation applies and $G$ is flat in $\log_2(1/\delta)$. For a predictor whose residual is its own error, that error is sized by the model rather than by $\delta$, so each halving costs roughly one extra bit per channel per step. The slope of $G$ therefore measures how much of a failure comes from the world being unpredictable rather than from the predictor being wrong, which is not what an accuracy measure reports. Two models can share a mean squared error while one loses to chaos and the other to its own mistakes.

\section{Experiments}

\paragraph{Setup.}
We evaluate four pretrained models from the Chronos Bolt family at $9$, $21$, $48$ and $205$ million parameters \citep{ansari2024chronos}, against persistence, which repeats the last observed value, and a delay embedding with local linear analog prediction \citep{takens1981, farmer1987}. Each is measured on eight held out trajectories of $4096$ samples at five precisions between $0.8$ and $0.05$, giving $150$ measurements. The released models are univariate, so channels are run independently and the corrections summed. All evaluation is inference only.

\begin{figure}[t]
  \centering
  \includegraphics[width=\textwidth]{fig1_precision_sweep.pdf}
  \caption{Information rate against inverse precision, one panel per system, ordered by entropy rate. The dotted line is the floor. The delay embedding descends as the precision tightens, which is the response of a predictor limited by the unpredictability of the system. Every pretrained model rises in a tight bundle with persistence, which is the response of a predictor limited by its own error.}
  \label{fig:sweep}
\end{figure}

\begin{table}[t]
  \centering
  \caption{Median excess information over the floor in nats per unit time, with the median precision slope and forward pass saving. Zero excess is optimal and a slope near zero indicates a predictor limited by the system rather than by itself. Per system slopes are in Appendix~\ref{app:slopes}.}
  \label{tab:main}
  \small
  \setlength{\tabcolsep}{4.5pt}
  \begin{tabular}{lrrrrrrr}
    \toprule
    Predictor & van der Pol & Rossler & Halvorsen & Lorenz & Chen & Slope & Passes \\
    Floor $h_{\mathrm{KS}}$ & 0.000 & 0.067 & 0.651 & 0.905 & 2.061 & & \\
    \midrule
    Delay embedding        & \phantom{0}0.00 & \phantom{0}0.98 & \phantom{00}9.73 & \phantom{00}9.21 & \phantom{0}27.74 & $-4.62$ & 12.6$\times$ \\
    Persistence            & 11.39 & 17.58 & \phantom{0}84.47 & \phantom{0}93.28 & 213.08 & 30.97 & 1.1$\times$ \\
    Chronos Bolt tiny (9M)    & 16.50 & 24.19 & 111.19 & 111.14 & 278.00 & 34.82 & 1.0$\times$ \\
    Chronos Bolt mini (21M)   & 12.58 & 19.35 & 105.04 & 110.66 & 251.22 & 32.92 & 1.1$\times$ \\
    Chronos Bolt small (48M)  & 13.73 & 20.25 & 104.98 & 102.16 & 245.17 & 31.78 & 1.1$\times$ \\
    Chronos Bolt base (205M)  & 13.19 & 20.47 & 102.20 & \phantom{0}96.61 & 235.47 & 29.76 & 1.1$\times$ \\
    \bottomrule
  \end{tabular}
\end{table}

\paragraph{Distance from the floor.}
The ordering in Table~\ref{tab:main} is identical on all five systems and does not depend on the entropy rate. The delay embedding is closest by roughly an order of magnitude, persistence comes next, and every pretrained model at every size costs more information than repeating the last observed value, by $93.3$ against $96.6$ nats per unit time on Lorenz and $213.1$ against $235.5$ on Chen.

\paragraph{The precision response.}
Figure~\ref{fig:sweep} shows the curves and the separation is qualitative rather than a matter of degree. The delay embedding is at or below zero on every system, reaching exactly zero on the limit cycle, while every pretrained model is clearly positive everywhere and sits inside the band set by persistence. On Chen, persistence is at $76.8$ and the four pretrained models are at $89.1$, $77.8$, $74.6$ and $69.8$. By the measure that distinguishes tracking dynamics from treating them as noise, these models are not separable from having no model at all.

\paragraph{Scaling.}
Across a twentyfold range in parameters the excess falls by $11.46$ nats per unit time for each tenfold increase, and extrapolating to zero gives $10^{16.8}$ parameters, with a per system range from $10^{15.8}$ to $10^{21.8}$. Read as an order of magnitude rather than a prediction, it says closing this gap by scale alone is not a plausible route. Within the range measured the improvement is small and not monotone, since the middle size is best on several systems. The forward pass column of Table~\ref{tab:main} carries the same message, with the pretrained models needing consultation at essentially every step while the delay embedding needs one in twelve.

\paragraph{The instrument reaches the floor.}
Van der Pol is the calibration case, since its floor is exactly zero and any excess belongs to the predictor. The delay embedding records $0.001$ nats per unit time there, flat across all five precisions, so the floor is attainable and the gaps elsewhere belong to the predictors rather than to the coder. It is also the easiest signal in the benchmark, and every pretrained model still spends between $12.6$ and $16.5$ nats per unit time on it.

\section{Conclusion}

Forecasting evaluation has no zero, and that is solvable on systems where the optimum follows from the equations. Our coder measures against that optimum and our sweep separates the unpredictability of the world from the error of the model. Applied to pretrained temporal foundation models it reports that they spend more information than a trivial baseline on every system tested, that their precision response is indistinguishable from having no internal model, and that scale closes the gap far too slowly to reach the floor. The framework applies to any predictor that can extend a context, with no labels required.

\nocite{*}
\bibliographystyle{plainnat}
\bibliography{references}

\newpage
\appendix

\section{Per system precision slopes}
\label{app:slopes}

\begin{table}[h]
  \centering
  \caption{Slope of the information rate against $\log_2(1/\delta)$, per system.}
  \small
  \begin{tabular}{lrrrrr}
    \toprule
    Predictor & van der Pol & Rossler & Halvorsen & Lorenz & Chen \\
    \midrule
    Delay embedding        & \phantom{-}0.00 & $-0.70$ & $-5.35$ & $-4.62$ & $-10.63$ \\
    Persistence            & 3.79 & 6.28 & 30.97 & 31.53 & 76.79 \\
    Chronos Bolt tiny (9M)   & 4.76 & 7.88 & 35.07 & 34.82 & 89.08 \\
    Chronos Bolt mini (21M)  & 3.55 & 6.44 & 32.92 & 33.23 & 77.76 \\
    Chronos Bolt small (48M) & 3.92 & 6.90 & 31.78 & 30.95 & 74.61 \\
    Chronos Bolt base (205M) & 4.11 & 7.16 & 29.76 & 28.58 & 69.84 \\
    \bottomrule
  \end{tabular}
\end{table}

\section{Scaling}

\begin{figure}[h]
  \centering
  \includegraphics[width=0.62\textwidth]{fig2_scaling.pdf}
  \caption{Excess information against parameter count, with the pooled fit extended to where the excess would reach zero. Per system the extrapolated size ranges from $10^{15.8}$ for Chen to $10^{21.8}$ for Halvorsen.}
  \label{fig:scaling}
\end{figure}

\begin{table}[h]
  \centering
  \caption{Response of excess information to model size, and the extrapolated size at which the excess would reach the floor.}
  \label{tab:scaling}
  \small
  \begin{tabular}{lrr}
    \toprule
    System & Excess reduction per $10\times$ parameters & Zero excess at \\
    \midrule
    van der Pol & $-1.70$ & $10^{16.2}$ \\
    Rossler     & $-2.18$ & $10^{17.7}$ \\
    Halvorsen   & $-7.95$ & $10^{21.8}$ \\
    Lorenz      & $-12.80$ & $10^{16.3}$ \\
    Chen        & $-32.66$ & $10^{15.8}$ \\
    \midrule
    Pooled      & $-11.46$ & $10^{16.8}$ \\
    \bottomrule
  \end{tabular}
\end{table}

\section{Systems and their invariants}

Table~\ref{tab:systems} gives the five systems, their dimension, the integration and sampling intervals, the largest exponent and the entropy rate. All systems are integrated with a fixed step fourth order Runge and Kutta scheme using the analytic Jacobian for the variational equation. Lyapunov spectra are computed by the standard algorithm with periodic reorthonormalisation over $60000$ steps after a transient of $10000$ steps \citep{benettin1980}. An autonomous flow on a non fixed point attractor has exactly one zero exponent, and finite time estimation leaves it slightly off zero, so it is pinned to zero before the positive exponents are summed.

\begin{table}[h]
  \centering
  \caption{Benchmark systems. Exponents are in nats per unit time.}
  \label{tab:systems}
  \small
  \begin{tabular}{lrrrrrr}
    \toprule
    System & Dim. & Integration step & Sample interval & $\lambda_{\max}$ & $h_{\mathrm{KS}}$ & Samples per $1/\lambda_{\max}$ \\
    \midrule
    van der Pol & 2 & 0.020 & 0.200 & 0.0000 & 0.0000 & n/a \\
    Rossler     & 3 & 0.020 & 0.140 & 0.0670 & 0.0670 & 106.6 \\
    Halvorsen   & 3 & 0.005 & 0.040 & 0.6511 & 0.6511 & 38.4 \\
    Lorenz      & 3 & 0.005 & 0.045 & 0.9054 & 0.9054 & 24.5 \\
    Chen        & 3 & 0.002 & 0.016 & 2.0612 & 2.0612 & 30.3 \\
    \bottomrule
  \end{tabular}
\end{table}

Parameters follow standard settings. Lorenz uses $\sigma = 10$, $\rho = 28$ and $\beta = 8/3$ \citep{lorenz1963}. Rossler uses $a = b = 0.2$ and $c = 5.7$ \citep{rossler1976}. Chen uses $a = 35$, $b = 3$ and $c = 28$ \citep{chen1999}. Halvorsen uses $a = 1.4$ and van der Pol uses $\mu = 2$ \citep{sprott2010, vanderpol1926}. Against published values our estimator gives $0.887$ for Lorenz and $0.067$ for Rossler, and returns exactly zero for the limit cycle and two positive exponents for hyperchaotic test systems, which are the checks we use to validate it.

\section{Complete measurements}

Tables~\ref{tab:full1} and \ref{tab:full2} give every measurement behind the results in the main text, namely the information rate in nats per unit time at each precision for each predictor and system. For van der Pol the floor is zero so the information rate and the excess coincide. All $150$ measurements satisfy the reconstruction bound $|x_t - r_t| \le \delta/2$ to numerical precision, which is the correctness check on the coder.

\begin{table}[h]
  \centering
  \caption{Information rate in nats per unit time, van der Pol, Rossler and Halvorsen.}
  \label{tab:full1}
  \small
  \begin{tabular}{llrrrrr}
    \toprule
    System & Predictor & $\delta = 0.8$ & $0.4$ & $0.2$ & $0.1$ & $0.05$ \\
    \midrule
    van der Pol & Delay embedding & 0.001 & 0.001 & 0.001 & 0.007 & 0.020 \\
    & Persistence & 5.977 & 8.060 & 11.390 & 15.374 & 21.290 \\
    & Chronos Bolt tiny & 8.411 & 12.173 & 16.496 & 21.658 & 27.483 \\
    & Chronos Bolt mini & 6.497 & 9.480 & 12.577 & 16.080 & 20.960 \\
    & Chronos Bolt small & 6.963 & 9.921 & 13.733 & 17.716 & 22.649 \\
    & Chronos Bolt base & 6.717 & 9.503 & 13.187 & 17.730 & 23.136 \\
    \midrule
    Rossler & Delay embedding & 3.891 & 2.215 & 0.885 & 0.921 & 1.046 \\
    & Persistence & 9.860 & 13.506 & 17.644 & 24.638 & 35.683 \\
    & Chronos Bolt tiny & 10.748 & 17.542 & 24.259 & 31.563 & 43.137 \\
    & Chronos Bolt mini & 8.965 & 13.418 & 19.422 & 25.971 & 34.877 \\
    & Chronos Bolt small & 8.643 & 13.792 & 20.321 & 26.881 & 36.614 \\
    & Chronos Bolt base & 8.463 & 13.747 & 20.541 & 27.578 & 37.341 \\
    \midrule
    Halvorsen & Delay embedding & 32.952 & 17.765 & 10.323 & 9.445 & 10.377 \\
    & Persistence & 46.761 & 63.016 & 85.119 & 122.100 & 172.058 \\
    & Chronos Bolt tiny & 56.338 & 83.260 & 111.841 & 149.735 & 198.457 \\
    & Chronos Bolt mini & 52.802 & 78.626 & 105.692 & 140.421 & 186.506 \\
    & Chronos Bolt small & 55.220 & 80.062 & 105.631 & 139.212 & 184.561 \\
    & Chronos Bolt base & 52.725 & 77.562 & 102.854 & 132.567 & 174.012 \\
    \bottomrule
  \end{tabular}
\end{table}

\begin{table}[h]
  \centering
  \caption{Information rate in nats per unit time, Lorenz and Chen.}
  \label{tab:full2}
  \small
  \begin{tabular}{llrrrrr}
    \toprule
    System & Predictor & $\delta = 0.8$ & $0.4$ & $0.2$ & $0.1$ & $0.05$ \\
    \midrule
    Lorenz & Delay embedding & 31.160 & 13.119 & 8.641 & 8.962 & 10.115 \\
    & Persistence & 51.854 & 69.425 & 94.183 & 131.689 & 178.396 \\
    & Chronos Bolt tiny & 55.577 & 81.549 & 112.046 & 150.402 & 195.231 \\
    & Chronos Bolt mini & 57.604 & 82.926 & 111.569 & 147.816 & 191.333 \\
    & Chronos Bolt small & 53.827 & 76.636 & 103.070 & 137.072 & 178.371 \\
    & Chronos Bolt base & 50.842 & 72.391 & 97.518 & 127.435 & 166.207 \\
    \midrule
    Chen & Delay embedding & 69.611 & 47.191 & 29.797 & 25.592 & 27.263 \\
    & Persistence & 113.590 & 157.067 & 215.143 & 304.312 & 423.917 \\
    & Chronos Bolt tiny & 134.815 & 204.078 & 280.062 & 375.583 & 494.471 \\
    & Chronos Bolt mini & 129.134 & 186.944 & 253.284 & 336.079 & 443.355 \\
    & Chronos Bolt small & 129.153 & 184.417 & 247.235 & 326.091 & 431.388 \\
    & Chronos Bolt base & 122.750 & 174.210 & 237.530 & 310.682 & 403.726 \\
    \bottomrule
  \end{tabular}
\end{table}

\begin{table}[h]
  \centering
  \caption{Median forward pass saving of event driven scheduling relative to a per step forecaster.}
  \label{tab:passes}
  \small
  \begin{tabular}{lrrrrr}
    \toprule
    Predictor & van der Pol & Rossler & Halvorsen & Lorenz & Chen \\
    \midrule
    Delay embedding & 32.00 & 18.04 & 12.46 & 12.06 & 10.79 \\
    Persistence & 1.44 & 1.15 & 1.04 & 1.01 & 1.10 \\
    Chronos Bolt tiny & 1.01 & 1.12 & 1.03 & 1.02 & 1.04 \\
    Chronos Bolt mini & 1.10 & 1.53 & 1.04 & 1.04 & 1.08 \\
    Chronos Bolt small & 1.10 & 1.38 & 1.04 & 1.04 & 1.09 \\
    Chronos Bolt base & 1.12 & 1.33 & 1.05 & 1.07 & 1.13 \\
    \bottomrule
  \end{tabular}
\end{table}

\section{Cost and quality under event driven scheduling}

Because the coder consults a predictor only when a correction fires, the forward pass count is itself a measure of how often the model is wrong, and Table~\ref{tab:passes} reports it. The delay embedding needs between $10.8$ and $32.0$ times fewer passes than a per step forecaster. The pretrained models sit between $1.03$ and $1.13$ times and so require consultation at essentially every step. Quality and running cost become the same quantity under this scheme, which is useful wherever a model must roll forward over long horizons.

\section{Reproducibility}

\paragraph{Data generation.}
Each system contributes twenty four training trajectories, used only to fit the delay embedding library, and eight evaluation trajectories, each of $4096$ samples after a transient of $2048$ integration steps. Initial conditions are drawn by Gaussian perturbation of a reference point at a relative scale of $0.05$. The sampling interval targets twenty five samples per Lyapunov time and is additionally bounded by one fortieth of the dominant spectral period, the latter estimated from the strongest non constant peak of the power spectrum. Trajectories are standardised per channel using training statistics, so a precision of $0.1$ means one tenth of a standard deviation on every system.

\paragraph{Predictors.}
Persistence repeats the last reconstructed value. The delay embedding uses embedding dimension one above the system dimension, a lag of three samples, twenty four neighbours, a library subsampled to twenty thousand points and a ridge term of $10^{-3}$ on the local linear solve. Neighbour search is exact. The pretrained models are the released Chronos Bolt checkpoints in the tiny, mini, small and base sizes, used through the public pipeline at a context length of $512$ with the median of the predictive distribution taken as the point forecast. No fine tuning of any kind is performed. Because these backbones are univariate, the cross channel structure they cannot exploit is part of the measured excess.

\paragraph{Coder settings.}
The forecast block is thirty two steps. Warm up is the larger of one hundred and twenty eight steps and the predictor's context length. Precisions are $0.8$, $0.4$, $0.2$, $0.1$ and $0.05$. The rate estimator is an adaptive context model of order three over the correction symbols with Krichevsky and Trofimov smoothing, with the context formed from the signs of the three preceding symbols.

\paragraph{Compute.}
The complete evaluation reported here runs is completed in roughly 2 hours on a single RTX-class GPU with twenty four gigabytes of memory. Generating the trajectories and their Lyapunov spectra takes an additional few minutes on a single processor core and its results are cached.

\paragraph{Scope.}
The floor exists in the form used here for deterministic systems with a computable spectrum. Extending the same measurement to observed data yields a relative comparison in interpretable units rather than an absolute one, since the sweep can still be run but the anchor is unavailable. The systems here are low dimensional and fully observed, and the pretrained models were trained predominantly on business and infrastructure series, so the measurement reports how these models behave on deterministic dynamics rather than on the distribution they were built for. All results are from a single data seed, with the spread across five systems and five precisions serving as the source of variation.

\section{Artifact}

The full implementation, including the system definitions with analytic Jacobians, the Lyapunov solver, the coder, the rate estimator, the predictor interface with wrappers for the pretrained checkpoints, and a test suite covering the numerical components, will be released upon acceptance. The test suite checks the Lyapunov estimator against published exponents, the rate estimator against the analytic entropy of a synthetic stream, the reconstruction bound of the coder and the forward pass saving of the event driven schedule.

\end{document}
